% 中文摘要與英文摘要（標題上下留白略縮，使中文摘要維持一頁）
\ctexset{chapter={beforeskip=0pt,afterskip=0.7cm}}
\chapter*{摘要}
\addcontentsline{toc}{chapter}{摘要}
\markboth{摘要}{摘要}

{\linespread{1.3}\selectfont\enlargethispage{\baselineskip}
現代超級電腦與 GPU 叢集越來越常在功耗上限下運作：機房依機櫃電力限電，學生超算競賽限制整個叢集的功率，NVIDIA GH200 也允許使用者設定整顆模組的功耗上限。在功耗上限下，GPU 會把時脈降到功耗剛好不超過上限，因此程式的速度取決於每次運算消耗的能量；而 GPU 的耗電又隨資料而變，位元翻轉越多越耗電。另一方面，許多科學應用實際需要的精度，遠少於 FP64 提供的 52 個尾數位元。

本專題研究「尾數歸零」：以軟體將運算元低位的尾數位元設為 0，不改變資料格式與演算法，使運算單元與資料線路的翻轉減少，再於功耗上限下把省下的功耗轉換為更高的時脈與吞吐量。我們在 14 節點的 GH200 叢集上，以成對交錯的量測與「等效功耗」（未修改的程式需要多給多少功耗上限才能追上）評估此做法，涵蓋矩陣乘法微基準、8 節點 HPL 與 HPL-MxP、Quantum ESPRESSO 與 PySCF 等科學軟體、第三方非矩陣乘法程式，以及 AI 訓練與推論。

結果顯示，在功耗受限且截斷直接實作於程式內時，所有受測、以矩陣乘法為主且具正確性檢查的完整程式皆獲得正面收益。8 節點 HPL 保留 36 位有效位數（含隱藏位，即 35 個尾數位元；殘差距門檻約百倍）在 400–800~W 快 4.0–11.0\%，相當於每節點多給 39–127~W 的功耗上限；HPL-MxP 保留 3 位快 1.8–12.7\%（對角優勢的測試矩陣）；Quantum ESPRESSO 依求解器的收斂門檻自動調整位元數，在總能量印出的每一位皆不變的情況下快 2.4–4.2\%、少耗電 2.5–3.9\%（金表面的收益約一半來自 Davidson 迭代減少）；DF-MP2 截斷到 20–32 位在 500–700~W 快 4.9–8.3\%；指數跨度寬且功耗高的高斯過程地統計截斷到 28 位在 500–700~W 快 7.4–8.2\%，估出的參數與原生相同，而 FP32 結果不可用、INT8 模擬較慢。在預設的 1000~W 上限下，HPL 仍快 3.9\%，每次求解少耗電 7.5\%。機制實驗顯示 62–97\% 的省電來自低位元不再翻轉，在 tensor core、FP64 乘加單元與記憶體搬移上皆成立。

適用範圍方面，與 INT8 模擬 FP64 的分界點隨矩陣形狀與功耗上限而變；所有矩陣乘法改用 FP32 或 3xTF32 時，長內積的 Quantum ESPRESSO 無法收斂；列內指數跨度寬時，固定位數的 INT8 模擬可能在沒有警告的情況下算錯（高斯基底 DFT 上 cuBLAS 固定 36 位模擬的能量誤差達 52\mEh，同樣保留 36 位的最近捨入截斷只差 $10^{-6}$\Eh），但跨度寬不一定導致錯誤，以 tile 為單位換格式也能處理這類資料，截斷還可疊加其上；在已使用 FP8 的 AI 推論上，FP8 已取得幾乎全部的好處。尾數歸零每個元素只多一次位元運算、不需要新硬體，最適合具殘差修正或自洽迭代的科學求解器。

{\small 註：本報告中所有截斷後的 HPL／HPL-MxP 與 MLPerf 數字皆為研究用途，並非 TOP500 或 MLPerf 的正式成績。\par}

\vspace{0.5em}
\noindent\textbf{關鍵詞：}功耗上限、能源效率、浮點精度、尾數截斷、GPU、混合精度、HPL
\par}

\ctexset{chapter={beforeskip=0.5cm,afterskip=1.2cm}}
\chapter*{Abstract}
\addcontentsline{toc}{chapter}{Abstract}
\markboth{Abstract}{Abstract}

{\linespread{1.3}\selectfont
We study a software knob that exploits data-dependent GPU power: zeroing unneeded low-order mantissa bits of FP64,
BF16 and FP8 operands. The study runs on 14 NVIDIA GH200 nodes under module power caps, with paired measurements and
an equivalent-power metric (the extra cap the unmodified program needs to catch up). Under a binding cap, every tested
GEMM-dominated HPC program that truncates internally and checks its results gains. At 400--800\,W, tuned 8-node HPL
keeping 36 significand bits (about $100\times$ residual margin) runs 4.0--11.0\% faster, worth 39--127\,W per node,
and HPL-MxP keeping 3 bits runs 1.8--12.7\% faster on a diagonally dominant test matrix. At 450--550\,W, Quantum
ESPRESSO, with the bit count tied to its Davidson tolerance, runs 2.4--4.2\% faster (partly from fewer iterations)
with every printed energy digit unchanged. At 500--700\,W, DF-MP2 keeping 20--32 bits runs 4.9--8.3\% faster, and a
wide-range Gaussian-process likelihood keeping 28 bits runs 7.4--8.2\% faster, where FP32 fails and INT8 emulation is
slower. At the default 1000\,W cap, HPL is still 3.9\% faster and uses 7.5\% less energy per solve. Fill controls
attribute 62--97\% of the saving to bits that stop toggling. Safe bit counts need validation across problem sizes,
and the crossover with INT8 emulation of FP64 moves with matrix shape and power cap. With wide per-row exponent
spans, fixed-width INT8 emulation can fail silently: at 36 bits it shifts a Gaussian-basis DFT energy by
52\,m$E_\mathrm{h}$, whereas rounding to 36 bits errs by $10^{-6}\,E_\mathrm{h}$. All-FP32 or 3xTF32 GEMMs fail on
long inner products, and in Hopper AI inference FP8 already captures most of the benefit. Truncated HPL, HPL-MxP and
MLPerf numbers in this report are research results, not official TOP500 or MLPerf scores.

\vspace{1em}
\noindent\textbf{Keywords:} power capping, energy efficiency, floating-point precision, mantissa truncation, GPU,
mixed precision, HPL
\par}
